\section{Introduction}\label{sec:introduction}

A projective symplectic variety carries a unique singular Ricci-flat
K\"ahler metric in each K\"ahler class, in the sense of
Eyssidieux--Guedj--Zeriahi \cite[Theorem~7.5 and Corollary~7.6]{egz09}.
The existence theorem does not by itself describe the local metric
geometry at a singular point. In a moduli space of bundles, a second
description is available: Hermitian--Einstein connections define a
hyperk\"ahler $L^2$ metric on the stable locus. We identify these two
metrics on a family of singular compact K3 moduli spaces and obtain
a quadratic comparison with explicit hyperk\"ahler cones at every
singular point.

The simplest example is already a compact space with two
positive-dimensional singular strata. Let $S$ be a quartic K3 surface
with $\operatorname{Pic}(S)=\ZZ H$, $H^2=4$, and put
\begin{equation}\label{eq:quartic}
 v_0=(5,H,0),\qquad v=2v_0=(10,2H,0),\qquad M=M_H(v).
\end{equation}
Here $M_H(v)$ denotes the moduli space of $H$-Gieseker-semistable
sheaves of Mukai vector $v$. This space has complex dimension $18$.
Every semistable sheaf is a vector bundle, and its slope
Jordan--H\"older factors are vector bundles as well. If
$N=M_H(v_0)$, then $N$ is a smooth compact hyperk\"ahler sixfold,
and the singular locus and its diagonal are
\begin{equation}\label{eq:quartic-strata}
 \Sigma=\Sym^2N,\qquad \Delta=\{[G\oplus G]:[G]\in N\}.
\end{equation}
Thus $\dim_{\CC}\Sigma=12$ and $\dim_{\CC}\Delta=6$.
The variety $M$ has terminal, locally factorial symplectic
singularities, admits no symplectic resolution, and has no smoothing.
These assertions follow from the moduli results recalled in
Section~\ref{sec:moduli}. The exclusion of a non-locally-free
boundary is proved there by a Mukai-pairing inequality.

\subsection{The family and the main result}
Let $S$ be a projective K3 surface with
$\operatorname{Pic}(S)=\ZZ H$, where $H$ is ample. Write the Mukai
vector of a sheaf as $v(F)=(\rk F,c_1(F),\chi(F)-\rk F)$, with
pairing
\[
 \langle(r,c,s),(r',c',s')\rangle=c\cdot c'-rs'-r's.
\]
We impose the conditions
\begin{equation}\label{eq:family}
 v_0=(r_0,aH,s_0),\qquad
 \gcd(r_0,aH^2)=1,\qquad q:=v_0^2\ge4,
 \qquad r_0\ge q+1,
\end{equation}
and set
\begin{equation}\label{eq:MN}
 M=M_H(2v_0),\qquad N=M_H(v_0).
\end{equation}
The integer $q$ is even. Fix the Ricci-flat K\"ahler form
$\omega\in c_1(H)$ on $S$.
Proposition~\ref{prop:no-bubbling} proves that slope and Gieseker
semistability agree for the sheaves under consideration.
The moduli results in Section~\ref{sec:moduli} identify
the stable locus $M^s$ with the regular locus of $M$. The
Hermitian--Einstein metrics therefore define the harmonic $L^2$
metric $g^{L^2}$ on $M_{\reg}$.

Here is the local cone that occurs in the theorem. For a polystable
bundle $E$, let
\begin{equation}\label{eq:cone-intro}
 V_E=H^1(S,\End E),\qquad K_E=\Aut(E,h_E),\qquad
 \Cone_E=V_E\hkr_0K_E.
\end{equation}
Harmonic representatives give $V_E$ its flat hyperk\"ahler
structure. The last expression means the simultaneous zero set of
the three moment maps, divided by the unitary automorphism group;
an ineffective scalar circle can be removed. The regular part of
this quotient carries the reduction metric $g_{\Cone_E}$. Its full
length metric is denoted by $d_{\Cone_E}$, and
$r(x)=d_{\Cone_E}(0,x)$. Equivalently, $\Cone_E$ is the flat
hyperk\"ahler quotient of the doubled Ext-quiver of the stable
summands of $E$. Section~\ref{sec:harmonic} specifies the harmonic
and moment-map conventions.

\begin{theorem}\label{thm:main}
Assume \eqref{eq:family}. There is an ample Donaldson determinant
line bundle $\lambda$ on $M$ and a constant $\kappa>0$ with the
following properties.
\begin{enumerate}
\item The K\"ahler form of $g^{L^2}$ extends to the canonical
singular Ricci-flat current in the class $\kappa c_1(\lambda)$.
It has continuous local potentials on $M$.
\item At every polystable point $[E]\in M$ there is a
biholomorphism of germs
\[
 \Psi_E:(\Cone_E,0)\longrightarrow(M,[E])
\]
for which, on the regular locus of a sufficiently small
neighbourhood,
\begin{equation}\label{eq:main-comparison}
 (1-C_Er^2)g_{\Cone_E}
 \le \Psi_E^*g^{L^2}
 \le (1+C_Er^2)g_{\Cone_E}.
\end{equation}
The same inequality holds on tangent vectors to every smooth
orbit-type stratum, with its harmonic metric. Within this fixed
chart, $C_E$ is independent of stabilizer type and of the ratios
of positive orbit eigenvalues.
\end{enumerate}
\end{theorem}

The determinant convention, including its positive sign, is given
in Section~\ref{sec:canonical}. The constants and chart radii in
\eqref{eq:main-comparison} may depend on the fixed bundle $E$,
its harmonic dimensions and its positive background elliptic
gaps. The theorem does not assert one radius or comparison
constant for all possible centred charts.

\begin{corollary}\label{cor:completion-main}
The metric completion of $(M_{\reg},g^{L^2})$ is a compact geodesic
space canonically homeomorphic to $M$. At every $p=[E]\in M$ its
metric tangent cone is unique and is $(\Cone_E,d_{\Cone_E},0)$.
In the chart of Theorem~\ref{thm:main},
\begin{equation}\label{eq:radius-main}
 d(p,\Psi_E(x))=r(x)(1+O_E(r(x)^2)).
\end{equation}
The Gromov--Hausdorff distance between the closed unit balls in
$(M,\epsilon^{-1}d,p)$ and in $\Cone_E$ is $O_E(\epsilon^2)$.
\end{corollary}

The completion assertion requires a path argument in addition to
the tensor estimate. We prove that paths along singular strata can
be approximated in length by regular paths. The perturbation is
transverse to the stratum in a product slice; radial dilation alone
preserves the singular stratum.

In the example \eqref{eq:quartic}, the two singular tangent cones
are
\begin{equation}\label{eq:quartic-cones}
 \begin{aligned}
 \Cone_{G_1\oplus G_2}
  &\cong \CC^{12}\times(T^*\CC^4\hkr_0\U(1)),
       &&G_1\not\cong G_2,\\
 \Cone_{G\oplus G}
  &\cong \CC^6\times
       ((\CC^6\otimes\mathfrak{sl}_2)\hkr_0\SU(2)).
 \end{aligned}
\end{equation}
The first normal factor is the closure of the minimal nilpotent
orbit in $\mathfrak{sl}_4$, of complex dimension six. The second
normal factor has complex dimension twelve. These products are
metric products, with the flat factors normalized by the harmonic
inner products.

The approximation in \eqref{eq:main-comparison} is not an exact
local isometry. Proposition~\ref{prop:not-isometric} proves that
no singular point has a neighbourhood isometric to a neighbourhood
of the vertex in its tangent cone. The reason is the intrinsic
geometry of the singular strata: their induced length metrics
contain the nonflat hyperk\"ahler metric of $N$. This does not
assert that the exponent two is optimal.

\subsection{Relation to earlier work}
The algebraic local models for these singularities belong to the
deformation theory of sheaves on K3 surfaces. Kaledin--Lehn--Sorger
established the structure and nonexistence of symplectic resolutions
for the relevant nonprimitive moduli spaces \cite{kls06}.
Arbarello--Sacc\`a identified the local quiver germs under
a formality hypothesis \cite[Propositions~4.4 and~6.1]{as18},
and Budur--Zhang proved that hypothesis for polystable sheaves
on projective K3 surfaces \cite[Theorem~1.2]{bz19}.
Perego--Rapagnetta established the irreducible
symplectic-variety properties used here \cite{pr23}.

On the analytic side, Itoh proved that the $L^2$ metric on the
irreducible Einstein--Hermitian moduli space over a K3 surface is
hyperk\"ahler \cite[Theorem~2.1 and Remark~(iii)]{itoh88}.
Verbitsky's theory of hyperholomorphic and projectively
hyperholomorphic bundles gives the underlying quadraticity
mechanism \cite[Theorems~6.2, 7.1 and 11.2, arXiv version]{verbitsky96}.
We use a surface calculation to show that the obstruction map is
quadratic in the particular harmonic coordinates needed for the
metric estimate. An unspecified analytic change of coordinates
would not preserve the stated quadratic error.

The quantitative local comparison for the harmonic metric is
imported from \cite[Theorem~1.1(iii)]{wuebben15}. Its application
here uses the full harmonic Kuranishi locus, including repeated
stable factors. The finite-dimensional metric slices and
localization arguments are those of \cite[Sections~4--6]{wuebben14}.
We state the required versions below. The additional global step
uses continuous Weil--Petersson potentials from
Buchdahl--Schumacher \cite[Theorem~11]{buchdahl21}, the curvature
formula of Biswas--Schumacher \cite[Proposition~4.3]{biswas16},
and uniqueness of the singular Ricci-flat current. Continuous
Weil--Petersson potentials themselves are an established input.

Namikawa--Odaka study canonical metrics and conical structures on
symplectic singularities, including a theorem under a symplectic
resolution or smoothing hypothesis and higher-order asymptotics
at isolated singularities \cite[Theorem~1.1 and Corollary~6.8]{no26}.
Odaka's subsequent result gives a hyperk\"ahler cone structure on
every symplectic singularity germ \cite[Corollary~1.3]{odaka26}.
The distinction relevant here is between the existence of a local
cone metric and the asymptotics of the metric selected by an ample
class on a compact moduli space. Theorem~\ref{thm:main} concerns
the latter metric and gives a relative estimate across its
positive-dimensional singular strata. It makes no assertion about
higher derivatives near lower strata or about the complete
Donaldson--Sun degeneration data.

\subsection{Proof and organization}
The local argument begins with the canonical harmonic
reconstruction $a=x-\bar\partial^*G(a\wedge a)$.
On a Ricci-flat K3 surface, harmonic endomorphism-valued
two-forms have parallel coefficients. Integration by parts
then annihilates the harmonic projection of every product
containing the coexact correction. Thus the obstruction
is exactly quadratic in the coordinates in which the
Hermitian--Einstein comparison has quadratic error.
The holomorphic family and the metric reconstruction
classify the same complex gauge orbit, so the resulting
chart is simultaneously holomorphic and quantitative.
Transverse perturbations in product slices then identify
the full stratified length distance with the completion
of the regular-locus metric.

For the global identification, the rank inequality forces
the defect cycles of all slope graded objects to vanish.
The Donaldson--Uhlenbeck morphism is consequently finite,
making its determinant polarization ample. A discriminant
calculation identifies this class with the harmonic
K\"ahler class. Continuous local potentials, together
with a depth-based cohomology restriction theorem, produce
a global continuous potential in that class. The
hyperk\"ahler volume identity then places the resulting
current in the uniqueness theorem of
Eyssidieux--Guedj--Zeriahi.

Section~\ref{sec:moduli} proves that the rank bound excludes both
non-locally-free sheaves and non-locally-free graded factors, and
describes the two singular strata. Section~\ref{sec:harmonic}
proves exact harmonic quadraticity and applies the quantitative
Hermitian--Einstein comparison in the same holomorphic quotient
coordinates. Section~\ref{sec:distances} identifies the completion
distance and derives the rescaling estimates.
Section~\ref{sec:canonical} identifies the ample determinant class
and proves the global Ricci-flat metric statement.
Section~\ref{sec:strata} proves the failure of local isometry.

All dimensions below are complex unless otherwise indicated.
We use $d^c=\frac{i}{2}(\bar\partial-\partial)$, so
$\ddc=i\partial\bar\partial$. A smooth stratum carries its
harmonic $L^2$ metric; a distance on the entire stratified space
will always be specified by its allowed paths.
